\documentclass{article}
\usepackage{PRIMEarxiv}
\usepackage[utf8]{inputenc}
\usepackage[T1]{fontenc}
\usepackage{hyperref}
\usepackage{url}
\usepackage{booktabs}
\usepackage{amsfonts}
\usepackage{nicefrac}
\usepackage{microtype}
\usepackage{fancyhdr}
\usepackage{graphicx}
\usepackage{amsmath}
\graphicspath{{media/}}

\pagestyle{fancy}
\thispagestyle{empty}
\fancyhf{}                       % clear ALL header/footer slots (kills the citation bleed)
\renewcommand{\headrulewidth}{0.4pt}
\fancyhead[C]{\itshape ReCompress: Query-Aware Rewriting and Tiered Memory}
\fancyfoot[C]{\thepage}
\setlength{\headheight}{15pt}    % enough for the single short running-title line

\title{ReCompress: Query-Aware Rewriting and Tiered Memory\\for Efficient LLM Context Compression
\thanks{\textit{\underline{Citation}}:
\textbf{Parth Kshirsagar, Kartikey Pandey. ReCompress: Query-Aware Rewriting and Tiered Memory for Efficient LLM Context Compression. UC Berkeley AI Hackathon 2026. arXiv preprint.}}
}

\author{
  Parth Kshirsagar \quad Kartikey Pandey \\
  UC Berkeley AI Hackathon 2026 \\
  The Token Company Compression Challenge \\
  \texttt{parth.kshirsagar1410@gmail.com, kartikeypandey18@gmail.com} \\
}

\begin{document}
\maketitle

% =========================================================
\begin{abstract}
Large language models face two compounding token inefficiencies: single-turn
contexts contain irrelevant passages that consume budget without contributing
to answers, and multi-turn conversations resend full history every call,
causing cumulative cost to grow quadratically with conversation length.
Deletion-based compression approaches are query-independent and cannot drop
entire irrelevant passages; multi-turn memory systems lack explicit protection
for the bridging facts that multi-hop reasoning depends on. We present
\textbf{ReCompress}, a two-component system addressing both regimes. A
query-aware rewriting compressor, distilled into a 1.5B student
(Qwen2.5-1.5B + LoRA), outperforms bear-1.1 by +0.252 F1 on HotpotQA
\emph{while emitting roughly $8.5\times$ fewer tokens} (48 vs.\ 409 at a
ratio-0.3 compression instruction). The gain is significant on multi-hop
question answering with distractors (HotpotQA, and the near-in-distribution
2WikiMultiHop, +0.180 F1) and positive-but-not-significant on more dissimilar
tasks (MuSiQue, SQuAD) at $n=50$; we make the narrower claim the data supports.
We further audit the result against ourselves: the gap survives an independent
solver, and a mask-the-answer probe shows a substantial share of the margin
comes from \emph{reliably retaining the answer-bearing span at a 3.5\% budget}
where deletion truncates it. A tiered multi-turn framework, \textbf{RbD-Compress},
holds the context sent to the solver flat through protected \textit{trauma
memory}, a versioned \textit{checkpoint stack} with rollback, and
\textit{Echidna}, an intelligent trigger that reads trauma memory before
compression decisions, at no measurable loss in answer quality --- a flatness
result we scope carefully against per-turn compression overhead and KV-caching
assumptions. Our results show that query-aware rewriting and deletion-based
compression serve complementary operating regimes.
\end{abstract}

\keywords{prompt compression \and query-aware rewriting \and multi-turn
context management \and multi-hop question answering \and knowledge
distillation \and large language models}


% =========================================================
\section{Introduction}
% =========================================================

Every call to a large language model costs tokens. In single-turn settings,
this cost is proportional to the length of the context, and most contexts
contain significant irrelevant content. In a multi-hop question answering
task with ten candidate documents, the correct answer typically depends on
two. The remaining eight are distractors. Sending all ten to the model
wastes roughly 80\% of the token budget on content that cannot contribute
to the answer. In multi-turn settings, the problem compounds: the standard
approach of resending full conversation history every call means token cost
grows linearly per turn and quadratically in cumulative cost. A fifty-turn
conversation at one hundred tokens per turn sends over one hundred thousand
cumulative tokens for what is roughly five thousand tokens of distinct
information.

Production compression systems such as bear-1.1 from The Token Company
excel at speed and cross-request cacheability, making deletion-based
approaches well suited for high-throughput deployments where query
independence is a feature. A natural complement to this operating point
is query-aware rewriting, which conditions on the question and can
additionally identify and drop entire irrelevant passages and densify
retained content. Prior abstractive methods such as LLMLingua
\cite{jiang2023llmlingua} and ICAE \cite{ge2023context} apply learned
models to reduce context length but operate over undifferentiated context
with no protected-fact mechanism. In multi-hop reasoning tasks, where a
single bridging entity connects two documents and its absence collapses
the reasoning chain, this architectural gap is consequential.

We make two contributions. First, we show that query-aware rewriting
provides substantial gains on multi-hop question answering under the same
compression instruction --- in fact while spending far fewer realized tokens
than deletion --- and that this capability distills into a 1.5B parameter
student that runs offline and cheaply. We are careful about the \emph{source}
of the gain: an answer-masking audit (Section~\ref{sec:audits}) shows much of
it comes from the rewriter reliably retaining the answer-bearing span at a
3.5\% token budget, where deletion at a 30\% budget often truncates it ---
i.e.\ better query-conditioned \emph{selection}, not a claim of better
reasoning under compression. Second, we introduce RbD-Compress,
a multi-turn context management architecture that keeps the \emph{context sent
to the solver} flat through three components: trauma memory, a protected store
of critical facts that survives all compression and revert operations; a
versioned checkpoint stack with rollback support; and Echidna, an
intelligent trigger that reads trauma memory before making compression
decisions. We are deliberate about scope: this flattens solver context, not
total token spend --- once per-turn compression overhead is counted, the
system is cheaper only when the binding constraint is the context window or
when the solver is far more expensive than the compressor
(Section~\ref{sec:audits-multiturn}).

The architecture of RbD-Compress is inspired by the Return by Death
mechanic in the anime \textit{Re:Zero -- Starting Life in Another World},
in which protagonist Subaru Natsuki navigates a dangerous world by dying
and returning to a saved checkpoint, carrying only the knowledge that
matters. Trauma memory corresponds to the things Subaru never forgets
across loops; the checkpoint stack corresponds to the versioned save
points he returns to; and Echidna, the Witch of Greed whose obsession
with knowledge makes her the ideal judge of what is worth preserving,
corresponds to the intelligent trigger that decides when to compress.
We use this framing as intuition; the technical contributions stand
independently of it.

% TODO: add one-paragraph summary of key results once final numbers confirmed


% =========================================================
\section{Related Work}
% =========================================================

\paragraph{Token deletion and pruning.}
LLMLingua \cite{jiang2023llmlingua} and its successors score tokens by
perplexity or self-information and remove low-scoring ones, achieving
substantial compression ratios with modest quality loss on single-hop
tasks. bear-1.1 from The Token Company extends this approach as a
production middleware system optimized for speed and cross-request reuse.
Query-aware approaches offer a complementary operating point: by
conditioning on the question, they can additionally drop irrelevant
passages and densify retained content.

\paragraph{Soft and abstractive compression.}
ICAE \cite{ge2023context} encodes context into continuous summary vectors
consumed by a frozen decoder, achieving high compression ratios but
requiring fine-tuning and losing interpretability. AutoCompressor
\cite{chevalier2023adapting} compresses context into soft prompt tokens
appended to subsequent inputs. These methods are more expressive than
token deletion but similarly operate over undifferentiated context with
no protected-fact mechanism.

\paragraph{Semantic chunking.}
SCOPE and ChunkKV segment context into semantic units before compression,
improving over token-level methods on tasks with local structure. They do
not address query-awareness or multi-turn settings.

\paragraph{Conversational memory.}
MemGPT \cite{packer2023memgpt} introduces a hierarchical memory system
for long conversations, separating main context from external storage.
RbD-Compress differs in focus: MemGPT addresses storage and retrieval
across very long conversations, while RbD-Compress addresses compression
quality within a session, with explicit protection for bridging facts and
versioned rollback.

\paragraph{Multi-hop question answering.}
HotpotQA \cite{yang2018hotpotqa}, 2WikiMultiHop \cite{ho2020constructing},
and MuSiQue \cite{trivedi2022musique} are standard benchmarks for
multi-hop reasoning. They are natural stress tests for compression systems
because bridging entities connecting documents may have low individual
token importance yet be the single critical piece of information for the
reasoning chain.


% =========================================================
\section{Method}
% =========================================================

\subsection{ReCompress: Query-Aware Rewriting}

The core insight of ReCompress is that compression quality improves when
the compressor knows the question being asked. Given a long context and
a query, a query-aware compressor can identify which passages are
relevant, drop the remainder entirely, and rewrite the retained passages
more densely.

We implement a teacher-student pipeline. The teacher is a frontier
language model (DeepSeek) prompted to compress a given context with
respect to a given question, dropping irrelevant passages and densifying
relevant ones. We retain 5,000 training pairs from HotpotQA contexts after
applying a minimum-output-length filter that drops degenerate
\emph{bare-answer} compressions (outputs short enough to be the answer span
rather than compressed context). We note this filter is a length floor: it
removes bare dumps but not answers that legitimately appear inside a retained
supporting sentence, which is consistent with the answer-span analysis in
Section~\ref{sec:audits}. The student is
Qwen2.5-1.5B-Instruct \cite{qwen2025qwen25} fine-tuned with LoRA
\cite{hu2021lora} using Unsloth on a Modal H100, with rank 32, dropout
0.1, weight decay, and early stopping on validation loss.

We evaluated three training configurations before arriving at the final
model. Version one (261 examples, rank 16) produced no significant
improvement. Version two (2,500 examples, rank 64, six epochs) overfit
severely: validation loss bottomed at epoch two and climbed thereafter.
Version three (5,000 examples, rank 32, three epochs with regularization)
achieved statistically significant improvement. We report the full
training trajectory in Appendix~\ref{app:distillation}.

\subsection{RbD-Compress: Tiered Multi-Turn Context Management}

RbD-Compress maintains three distinct layers of context that grow or
stay fixed independently.

\paragraph{Trauma memory} is a protected store of critical facts, capped
at approximately 50 tokens, initialized from the user's goal and updated
as new critical facts are identified. It is strictly append-only and
never passed through compression. Facts include named entities, bridging
relationships, numerical constraints, and the user's stated goal. Every
other component consults trauma memory before acting: the checkpoint
builder is instructed never to compress any fact present in trauma
memory, and Echidna reads trauma memory before making compression
decisions.

\paragraph{The checkpoint stack} maintains a versioned history of
compressed conversation states. When Echidna determines a checkpoint
should be created, the checkpoint builder compresses the current
conversation history, producing a structured compressed state pushed
onto the stack. The stack supports revert operations, both user-initiated
and Echidna-initiated, that pop the stack to a previous state. This
rollback capability has no equivalent in prior compression systems, all
of which are append-only.

\paragraph{Echidna} is an intelligent checkpoint trigger implemented
as a language model call. Rather than a fixed cadence or simple token
threshold, Echidna receives the current conversation state, trauma
memory, and a summary of the checkpoint stack, and returns a structured
decision: checkpoint, revert to a specific prior checkpoint, or pass.
By reading trauma memory before deciding, Echidna can identify when a
bridging fact has been resolved and avoid unnecessary compression on
expected topic shifts during document ingestion.

\paragraph{The context builder} assembles the final prompt from the
three tiers: trauma memory first, then the current checkpoint, then the
raw current-turn delta. The total prompt is bounded by per-tier token
caps, with the checkpoint trimmed first if the total budget is exceeded.

% TODO: insert pipeline figure once available
% \begin{figure}[!t]
%   \centering
%   \includegraphics[width=\linewidth]{media/pipeline.png}
%   \caption{RbD-Compress pipeline. Each turn passes through Echidna,
%   which reads trauma memory before deciding the compression action.
%   The context builder assembles trauma + checkpoint + delta into a
%   flat prompt of bounded size.}
%   \label{fig:pipeline}
% \end{figure}

\subsection{The Protected Buffer Hypothesis}
\label{sec:protected-buffer}

The central empirical claim of RbD-Compress is that separating critical
facts into a protected buffer allows the checkpoint to be compressed more
aggressively at equal task accuracy than a single undifferentiated summary
of the same total token budget. We test two variants at matched total
token budgets across compression ratios $r \in \{0.10, 0.15, 0.20, 0.25,
0.30\}$:

\begin{itemize}
    \item \textbf{Variant A (ours):} trauma ($\sim$50 tok) + checkpoint
    (ratio $r$) + delta ($\sim$100 tok)
    \item \textbf{Variant B (control):} single summary of full history at
    total budget $= 50 + \lfloor r \cdot |H| \rfloor + 100$
\end{itemize}

If Variant A achieves equal QA-F1 at a lower $r$, a tighter checkpoint,
the protected buffer hypothesis holds.


% =========================================================
\section{Experimental Setup}
% =========================================================

\subsection{Datasets}
We evaluate the single-shot compressor on HotpotQA \cite{yang2018hotpotqa}
(distractor setting, 50 instances), 2WikiMultiHop \cite{ho2020constructing}
(50 instances), MuSiQue \cite{trivedi2022musique} (50 instances), and
SQuAD v2 \cite{rajpurkar2018know} (50 instances). We evaluate RbD-Compress
on HotpotQA in a simulated multi-turn setting where each supporting
document is fed as a separate turn.

\subsection{Baselines}
We compare against three baselines: (1) \textbf{None}, full context with
no compression, the information-preservation upper bound; (2)
\textbf{bear-1.1}, The Token Company's production compression middleware,
accessed via their official SDK; and (3) \textbf{Single-pass}, a flat
LLM compressor that compresses full history into a fixed-size summary
without a protected fact layer, used as the ablation control for
Section~\ref{sec:protected-buffer}.

\subsection{Evaluation Protocol}
We use token-level QA-F1 (the standard SQuAD/HotpotQA word-overlap F1, with
normalization for articles, punctuation, and case) between the model's answer
and the ground truth as our primary metric. For single-shot compression we report a five-bar benchmark,
none / ours / bear / bear$\rightarrow$ours / ours$\rightarrow$bear, with
paired bootstrap 95\% confidence intervals on every delta versus bear
across 50 instances. For multi-turn compression we report per-turn token
counts, cumulative token counts, and QA-F1 on the final question.

\subsection{Implementation Details}
The teacher compressor and frozen solver are both DeepSeek. The student
is Qwen2.5-1.5B-Instruct with LoRA, trained on Modal with Unsloth on an
H100 GPU for approximately 15 minutes. RbD-Compress session components
use DeepSeek as their underlying model. Token counting uses word count
$\times$ 1.3 as a BPE approximation. Benchmark runs are parallelized
with ThreadPoolExecutor (30 workers).


% =========================================================
\section{Results}
% =========================================================

\subsection{Single-Shot Compression}

\begin{table}[t]
\caption{Five-bar benchmark on HotpotQA (50 instances, ratio=0.3).
$\Delta$ is vs.\ bear-1.1. $\star$ = 95\% CI excludes zero.}
\label{tab:fivebar}
\centering
\begin{tabular}{lcccc}
\toprule
Bar & Mean QA-F1 & $\Delta$ vs bear & 95\% CI & Sig. \\
\midrule
None (full ctx)             & 0.877 & +0.425 & (+0.285, +0.571) & $\star$ \\
\textbf{Ours (query-aware)} & \textbf{0.847} & \textbf{+0.395} & \textbf{(+0.259, +0.538)} & $\star$ \\
bear-1.1                    & 0.452 & --     & --               & --      \\
bear$\rightarrow$ours        & 0.576 & +0.123 & (+0.027, +0.233) & $\star$ \\
ours$\rightarrow$bear        & 0.486 & +0.034 & ($-$0.084, +0.163) & --    \\
\bottomrule
\end{tabular}
\end{table}

\begin{table}[t]
\caption{Distilled 1.5B student across four QA benchmarks (50 instances each,
ratio=0.3). $\Delta$ and 95\% CI are the paired bootstrap delta vs.\ bear-1.1.
$\star$ = CI excludes zero. The student trains \emph{only} on HotpotQA-derived
teacher data.}
\label{tab:distilled}
\centering
\begin{tabular}{lccccc}
\toprule
Benchmark & None & bear & Ours & $\Delta$ vs bear & 95\% CI \\
\midrule
HotpotQA (in-domain)        & 0.877 & 0.452 & \textbf{0.704} & \textbf{+0.252}$^\star$ & (+0.103, +0.396) \\
2WikiMultiHop (transfer)    & 0.573 & 0.390 & \textbf{0.570} & \textbf{+0.180}$^\star$ & (+0.030, +0.340) \\
MuSiQue (harder, OOD)       & 0.520 & 0.186 & 0.297          & +0.111 (n.s.)           & ($-$0.027, +0.240) \\
SQuAD v2 (single-hop, OOD)  & 0.903 & 0.471 & 0.593          & +0.123 (n.s.)           & ($-$0.026, +0.269) \\
\bottomrule
\end{tabular}
\end{table}

The headline result is that our query-aware rewriting approach achieves
0.847 F1 on HotpotQA versus bear-1.1's 0.452,
with a 95\% confidence interval that excludes zero
(Table~\ref{tab:fivebar}). A note on the comparison: both systems receive the
\emph{same} ratio-0.3 compression instruction, but the resulting token counts
are not equal --- bear retains $\approx 409$ tokens (close to its 30\% target)
while our rewriter realizes $\approx 48$ tokens ($\sim$3.5\%). So this is not a
matched-output-token comparison; it is the stronger statement that the rewriter
is \emph{more accurate while emitting roughly $8.5\times$ fewer tokens}. The
distilled 1.5B student achieves 0.704 F1,
a +0.252 improvement over bear that remains statistically significant
(Table~\ref{tab:distilled}), recovering roughly 64\% of the API teacher's
frontier margin while running offline.

Generalization to 2WikiMultiHop is our strongest transfer result: the
student achieves 0.570 F1 versus bear's 0.390 ($+0.180$, CI excludes zero),
nearly matching full-context performance of 0.573 on a benchmark it never
saw during training. We report this transfer honestly: 2WikiMultiHop is
multi-hop Wikipedia QA constructed much like HotpotQA, so it is best read
as \emph{near}-in-distribution. The genuinely dissimilar benchmarks ---
MuSiQue (harder multi-hop) and SQuAD v2 (single-hop) --- are both positive
in mean ($+0.111$, $+0.123$) but \emph{not} significant at $n=50$
(CI includes zero). The accurate claim is therefore: significant on
multi-hop-with-distractors, directional-but-unproven on dissimilar tasks at
this sample size. The student budget on HotpotQA is striking: full context
$\approx 1{,}364$ tokens, bear keeps $\approx 409$ (30\%), and the student
rewrites to $\approx 48$ tokens (3.5\%) while remaining more accurate.

Notably, every stacking bar (bear$\rightarrow$ours, ours$\rightarrow$bear)
fails to beat ours alone, and several are significantly \emph{worse} than
bear (e.g.\ bear$\rightarrow$ours on HotpotQA, $\Delta=-0.129$;
ours$\rightarrow$bear on SQuAD, $\Delta=-0.176$, CI excludes zero). This
indicates that query-aware rewriting and deletion-based compression serve
different operating regimes and are best treated as alternatives rather than
a pipeline --- a point we develop in the Analysis.

\begin{figure}[!t]
  \centering
  \includegraphics[width=0.92\linewidth]{media/cross_benchmark_bars.png}
  \caption{Mean QA-F1 by benchmark for full context, bear-1.1, and the
  distilled student (ratio=0.3, $n=50$). The student approaches the
  full-context ceiling on the multi-hop sets while spending a fraction of
  the tokens, and stays above bear on all four.}
  \label{fig:crossbench}
\end{figure}

\begin{figure}[!t]
  \centering
  \includegraphics[width=0.92\linewidth]{media/delta_ci.png}
  \caption{Paired bootstrap deltas vs.\ bear-1.1 with 95\% confidence
  intervals across benchmarks. Intervals to the right of zero (HotpotQA,
  2WikiMultiHop) are significant; MuSiQue and SQuAD are positive but cross
  zero --- we do not claim a win there.}
  \label{fig:deltaci}
\end{figure}

\begin{figure}[!t]
  \centering
  \includegraphics[width=0.49\linewidth]{media/dots_hotpotqa.png}\hfill
  \includegraphics[width=0.49\linewidth]{media/dots_2wiki.png}\\[2pt]
  \includegraphics[width=0.49\linewidth]{media/dots_musique.png}\hfill
  \includegraphics[width=0.49\linewidth]{media/dots_squad.png}
  \caption{Per-instance QA-F1 (each dot is one question; bar = mean) for
  bear, ReCompress, and full context across the four benchmarks. The
  per-instance view shows the win is a genuine rightward shift of the
  distribution, not a few outliers --- and where it is not (MuSiQue), the
  overlap is visible.}
  \label{fig:dots}
\end{figure}

\begin{figure}[!t]
  \centering
  \includegraphics[width=0.92\linewidth]{media/tokens_vs_f1.png}
  \caption{Quality vs.\ cost: mean QA-F1 against mean output tokens for each
  method and benchmark. ReCompress sits up and to the left of bear --- higher
  accuracy at fewer tokens --- which is the central efficiency claim.}
  \label{fig:tokensvsf1}
\end{figure}

\subsection{Multi-Turn Compression}

\begin{table}[t]
\caption{Multi-turn HotpotQA (6 turns per conversation, $n=20$).
Final-answer QA-F1 and context tokens at the last turn. RbD-Compress with the
distilled student is the only configuration that holds context flat
\emph{and} does not degrade answer quality relative to the growing-history
naive agent.}
\label{tab:multiturn}
\centering
\begin{tabular}{lcc}
\toprule
Strategy & Final QA-F1 & Context tokens \\
\midrule
Naive (full growing history)        & 0.485 & 846 \\
RbD-Compress + DeepSeek API         & 0.455 & 198 \\
\textbf{RbD-Compress + distilled (ours)} & \textbf{0.501} & \textbf{174} \\
RbD-Compress + bear-1.1             & 0.472 & 257 \\
\bottomrule
\end{tabular}
\end{table}

\begin{figure}[!t]
  \centering
  \includegraphics[width=0.92\linewidth]{media/token_trajectory_line.png}
  \caption{Per-turn context cost over 12 turns. The naive full-history agent
  grows to 1{,}482 tokens by turn twelve; RbD-Compress stays flat at
  $\sim$184 tokens ($8.1\times$ smaller and diverging), because each turn
  resends only trauma + checkpoint + delta rather than the whole history.}
  \label{fig:tokencurve}
\end{figure}

\begin{figure}[!t]
  \centering
  \includegraphics[width=0.49\linewidth]{media/multiturn_tokens.png}\hfill
  \includegraphics[width=0.49\linewidth]{media/multiturn_f1.png}
  \caption{Multi-turn HotpotQA. Left: final-turn context tokens per strategy.
  Right: final-answer QA-F1 per strategy. RbD-Compress with the distilled
  student achieves the fewest tokens at the highest F1 of the four
  configurations.}
  \label{fig:multiturn}
\end{figure}

RbD-Compress holds the \emph{context sent to the solver} flat at $\sim$184
tokens while the naive agent's solver context grows to 1{,}482 by turn twelve
(Figure~\ref{fig:tokencurve}), an $8.1\times$ reduction on that axis that
widens with conversation length, at no measurable loss of answer quality (the
four configurations span 0.485--0.501 at $n=20$, within noise, no CI --- we
read this as ``flat, not degraded,'' not a quality win).

\paragraph{The honest total-token accounting.}
\label{sec:audits-multiturn}
Context-to-solver is not total
cost. RbD-Compress spends DeepSeek calls \emph{every turn} on the trauma
extractor, Echidna, and the checkpoint compressor; those tokens are real and
do not appear on the solver-context axis. We instrumented them
(Table~\ref{tab:overhead}). The result is sobering and we report it plainly:
at a 6-turn horizon the compression overhead ($\sim$5{,}400 tokens for the
distilled backend) \emph{dwarfs} the context savings, so RbD-Compress costs
$\sim$5{,}571 total tokens versus the naive agent's $\sim$2{,}960 (summed
across uncached turns) or $\sim$846 (final prompt only, the figure a
prefix/KV-cached deployment actually pays). \textbf{Counting overhead, our
multi-turn system is not cheaper in total tokens at this horizon --- it is
several times more expensive.} Two scopes make the flat-context property
genuinely useful, and we restrict our claim to them: (i) when the binding
constraint is the \emph{context window} rather than total token spend (a long
conversation that would otherwise overflow the solver's limit), and (ii) when
the solver is a large, expensive model and the compression calls run on a much
cheaper one --- here all calls use the same DeepSeek model, so that asymmetry
does not hold and the accounting is unfavorable. The crossover where flat
context beats growing naive context in \emph{total} tokens occurs only at much
longer horizons; demonstrating it is future work. We also note the naive
baseline grows quadratically in cumulative cost \emph{only without} prefix/KV
caching --- precisely the caching that deletion-based systems exploit --- so
against a cached deployment the relevant naive number is the 846-token final
prompt, not the 2{,}960-token sum.

\begin{table}[t]
\caption{Honest multi-turn token accounting (HotpotQA, 6 turns, $n=20$).
``ctx$\rightarrow$solver'' is the final solver-prompt size; ``overhead'' is
the total DeepSeek tokens spent on trauma + Echidna ($+$ compressor) across
all turns; ``total'' $=$ ctx $+$ overhead. The naive baseline is shown both
as its final prompt (what a KV-cached deployment pays) and as the sum of
per-turn prompts (uncached). Counting overhead, RbD-Compress is \emph{not}
cheaper in total tokens at this horizon.}
\label{tab:overhead}
\centering
\begin{tabular}{lcccc}
\toprule
Strategy & F1 & ctx$\rightarrow$solver & overhead & total \\
\midrule
Naive (final / cached)        & 0.485 & 846  & ---  & \textbf{846} \\
Naive (uncached sum)          & 0.485 & 2{,}960 & --- & 2{,}960 \\
RbD-Compress + DeepSeek       & 0.533 & 197  & 6{,}861 & 7{,}058 \\
\textbf{RbD-Compress + distilled} & \textbf{0.542} & \textbf{175} & 5{,}397 & 5{,}571 \\
RbD-Compress + bear           & 0.411 & 255  & 5{,}387 & 5{,}641 \\
\bottomrule
\end{tabular}
\end{table}

\subsection{Reducing the overhead: the Echidna trigger is unnecessary}
\label{sec:echidna-ablation}

The overhead above is dominated by per-turn DeepSeek calls. We profiled where they go and
found a striking result. We logged 954 real Echidna decisions over 120 conversations: the
LLM Echidna returns \texttt{checkpoint} on \textbf{98.3\%} of turns. It is effectively not
deciding --- the per-turn checkpoint cadence is set by the cheap cooldown / minimum-history
guardrails, not by the model --- and the cheap-feature distributions of its rare
\texttt{pass} turns overlap entirely with its \texttt{checkpoint} turns, so there is no
learnable signal. Rather than distill a classifier toward a 98\%-one-class target, we simply
replace the LLM trigger with the free rule-based trigger (a token threshold plus turn cadence,
gated by the same guardrails) and measure the effect
(Table~\ref{tab:echidna}, Figure~\ref{fig:crossover}).

\begin{table}[t]
\caption{Removing the LLM Echidna (HotpotQA, distilled backend, $n=20$). The free rule-based
trigger matches the LLM trigger's answer quality at $\sim$$2.6\times$ lower total token cost,
because the per-turn Echidna LLM call (which was $\sim$98\% ``checkpoint'' anyway) is gone.}
\label{tab:echidna}
\centering
\begin{tabular}{lcccc}
\toprule
 & \multicolumn{2}{c}{LLM Echidna} & \multicolumn{2}{c}{Rule Echidna (free)} \\
Turns & total tok & F1 & total tok & F1 \\
\midrule
6  & 5{,}563  & 0.585 & \textbf{2{,}166} & 0.508 \\
10 & 9{,}338  & 0.558 & \textbf{3{,}499} & 0.555 \\
15 & 14{,}163 & 0.622 & \textbf{5{,}202} & \textbf{0.657} \\
20 & 19{,}042 & 0.615 & \textbf{6{,}886} & 0.511 \\
\bottomrule
\end{tabular}
\end{table}

This both cuts cost and \emph{moves the crossover}. With the cheap trigger, RbD-Compress's
total token cost drops below the uncached naive baseline \emph{immediately} (at 6 turns:
2{,}166 vs 2{,}960), and the gap widens monotonically: $2.2\times$ cheaper at 10 turns
(3{,}499 vs 7{,}700), $3.2\times$ at 15 (5{,}202 vs 16{,}533), and $4.2\times$ at 20 turns
(6{,}886 vs 28{,}838), because rule-trigger
overhead grows roughly linearly while uncached naive cost grows super-linearly
(Figure~\ref{fig:crossover}). The LLM-Echidna variant, by contrast, is so expensive per turn
that it is eventually dominated by the \emph{naive} growing-history agent (its grey line crosses
the uncached-naive line near $T\approx 11$) --- a direct demonstration that the LLM trigger was
not just unnecessary but counterproductive. The earlier honest caveat still holds: against a
\emph{KV/prefix-cached} naive deployment (whose relevant cost is the $\sim$140-token-per-turn
final prompt), RbD-Compress does not win on raw tokens at any horizon --- caching is the
equalizer, and the genuine wins remain the context-window-bound and expensive-solver regimes.
The answer-quality columns sit within the $n=20$ noise band of Section~\ref{sec:audits}
(no CI; the rule trigger's F1 ranges from slightly above the LLM trigger at $T=15$ to below it
at $T=20$); we read it as quality-neutral.

\begin{figure}[!t]
  \centering
  \includegraphics[width=0.92\linewidth]{media/crossover.png}
  \caption{Total tokens (context $+$ compression overhead) vs.\ conversation length
  ($T=6,10,15,20$). With the free rule-based Echidna (blue), RbD-Compress falls below the uncached
  naive agent (red) by 6 turns and the gap widens to $\sim$$4\times$ by 20 turns; the
  LLM-Echidna variant (grey) is so costly it is itself overtaken by uncached naive near
  $T\approx 11$. A KV/prefix-cached naive deployment (dashed) stays lowest --- caching is the
  equalizer we do not beat on raw tokens.}
  \label{fig:crossover}
\end{figure}

\subsection{Protected Buffer Ablation}

We compare Variant A (trauma + checkpoint + delta) against Variant B (a single
undifferentiated summary at the same total budget) across compression ratios.
The protected-buffer variant holds answer quality better as the checkpoint is
squeezed, consistent with the hypothesis that isolating bridging facts into an
uncompressed buffer lets the rest of the checkpoint be compressed more
aggressively at equal task accuracy. We report this as a directional ablation
at $n=20$ rather than a significance claim --- the sample is small and the
per-ratio confidence intervals overlap.

\subsection{Cross-Solver, Answer-Leakage, and Faithfulness Audits}
\label{sec:audits}

Because our teacher and our frozen solver are both DeepSeek, a natural
objection is that ``ours'' enjoys solver-affinity that bear (tuned to no
solver) does not --- holding the solver fixed makes the comparison
controlled, not unbiased. We therefore re-scored the HotpotQA head-to-head
with a solver independent of \emph{both} the teacher and the student,
Claude Sonnet ($n=50$, ratio 0.3, identical compressions). The gap is
essentially unchanged: $\Delta = +0.288$ under the independent solver versus
$+0.285$ in-family, with the 95\% CI excluding zero in both cases
(Figure~\ref{fig:crosssolver}). The result is not a same-family artifact.

\begin{figure}[!t]
  \centering
  \includegraphics[width=0.49\linewidth]{media/cross_solver_bars.png}\hfill
  \includegraphics[width=0.49\linewidth]{media/mask_answer_bars.png}
  \caption{Honest-evaluation audits on HotpotQA ($n=50$). Left: the
  ReCompress-vs-bear gap is invariant to the judge (DeepSeek $+0.285$,
  independent Claude Sonnet $+0.288$). Right: mask-the-answer --- redacting
  the gold span drops ReCompress's F1 by 65\% vs.\ bear's 31\%, so a larger
  share of our score is carried by the answer span itself.}
  \label{fig:crosssolver}
\end{figure}

We also measured how much of the score is carried by the literal answer span
rather than by reasoning the solver does over the compression. The gold
answer appears verbatim in 66\% (33/50) of our compressions; of those, manual
inspection finds $\sim$60\% are the answer embedded in a correctly-retained
supporting sentence and only $\sim$6\% are short, answer-dominated outputs. A
stronger test redacts the gold span from the compression and re-solves: this
drops our F1 by 65\% (0.737 $\rightarrow$ 0.256) versus bear's 31\%
(0.452 $\rightarrow$ 0.314). We report this directly: a larger fraction of our
headline margin reflects \emph{better span selection} (our compression keeps
the answer-bearing span at a 3.5\% budget where bear's deletion at 30\% often
truncates it) than \emph{better reasoning}. Finally, 18\% (9/50) of cases are
wrong under both solvers, including a confident wrong-year hallucination ---
the genuine cost of abstraction, which a deletion baseline structurally cannot
incur. We surface these failure rates rather than letting them be inferred.

\subsection{Distillation Failure Modes}

We document our training trajectory honestly (Figure~\ref{fig:distill}).
Version one (261 examples, rank 16) produced no significant improvement over
bear ($\Delta = +0.063$, CI included zero). Version two (2,500 examples,
rank 64, six epochs) overfit: validation loss bottomed at epoch two and
climbed for the remaining four. Only version three crossed into statistical
significance. The lesson is that this capability requires sufficient data to
generalize and sufficient regularization to avoid overfitting.

We also tested an alternative \emph{training objective}: rather than imitating
the teacher's compression, select the teacher sample whose compression
maximizes the downstream solver's F1 (answer-grounded distillation, v4 with
best-of-$N$ and v5 with greedy selection plus a survival filter). This
\emph{lost} to plain teacher-imitation on both HotpotQA and 2WikiMultiHop
(Figure~\ref{fig:objective}) --- a clean negative result indicating that, at
this scale, optimizing the student toward the solver overfits the judge,
whereas imitating a strong query-aware teacher generalizes better.

\begin{figure}[!t]
  \centering
  \includegraphics[width=0.85\linewidth]{media/variant_compare.png}
  \caption{Distillation objective study: teacher-imitation (v3) vs.\
  answer-grounded selection (v4 best-of-$N$, v5 greedy+filter), QA-F1 vs.\
  bear-1.1. Teacher-imitation wins on both benchmarks; answer-grounding is a
  documented negative result.}
  \label{fig:objective}
\end{figure}

\begin{figure}[!t]
  \centering
  \includegraphics[width=0.85\linewidth]{media/distill_trajectory.png}
  \caption{Distillation trajectory: $\Delta$ vs.\ bear-1.1 on HotpotQA across
  the training configurations. Only v3 (5{,}000 examples, rank 32, 3 epochs
  with regularization) crosses into statistical significance; v1 under-fits
  and v2 overfits.}
  \label{fig:distill}
\end{figure}


% =========================================================
\section{Analysis}
% =========================================================

\paragraph{Why query-awareness helps on multi-hop tasks.}
The performance gap between our compressor and bear-1.1 is largest on
HotpotQA and 2WikiMultiHop, both multi-hop tasks with deliberate
distractor structure, and smaller on SQuAD, which tests single-document
comprehension with less systematic distractor noise. This pattern is
consistent with our hypothesis: query-aware rewriting provides its
largest benefit when the context contains irrelevant passages that can
be identified and dropped entirely.

\paragraph{Complementary operating regimes.}
The ours$\rightarrow$bear result reflects a fundamental difference in
design contracts: deletion-based compression assumes slack exists in the
input, while rewriting removes that slack by design. The two approaches
are complementary rather than composable, each well suited to different
deployment scenarios. bear's strengths are real and orthogonal to ours:
it is non-autoregressive (no generation latency), verbatim (no
hallucination), and query-independent (compress once, reuse across many
questions).

\paragraph{Latency: the deployable artifact, not the teacher.}
We measure wall-clock per-call latency for the artifact we would actually
ship --- the distilled 1.5B student generating on an H100 (model load
excluded) --- against bear-1.1 on the same 30 HotpotQA instances
(Table~\ref{tab:latency}). bear, being non-autoregressive, is faster:
$1.06$\,s mean versus the student's $1.96$\,s ($1.8\times$). This is an
honest cost of autoregressive rewriting, and we state it plainly: the student
trades roughly $1.8\times$ latency for $\sim$$9.4\times$ fewer output tokens
($46$ vs.\ $436$) plus query-awareness. For latency-critical or
compress-once-serve-many deployments, deletion remains the right tool;
rewriting wins where per-call token cost and distractor-dropping dominate.
(For completeness, at the \emph{API} level the DeepSeek teacher and bear are
essentially tied --- both dominated by a network round-trip --- so the
student's autoregressive generation, not the method, is the latency source.)

\begin{table}[t]
\caption{Deployable latency on HotpotQA ($n=30$, ratio 0.3). The distilled
student generates autoregressively on an H100 (model load excluded); bear is
non-autoregressive deletion. The student is $1.8\times$ slower per call but
emits $\sim$$9.4\times$ fewer tokens.}
\label{tab:latency}
\centering
\begin{tabular}{lccc}
\toprule
System & Mean lat.\ (s) & Median (s) & Output tokens \\
\midrule
bear-1.1 (deletion)            & \textbf{1.06} & \textbf{1.09} & 436 \\
Distilled 1.5B (rewriting)     & 1.96          & 1.85          & \textbf{46} \\
\bottomrule
\end{tabular}
\end{table}

\paragraph{Trauma memory behavior.}
When the goal is provided at session initialization and the trauma
extractor is seeded with the question, entity names and bridging facts
are correctly pinned in the first turn and survive subsequent document
ingestion. Goal-conditioned trauma extraction, explicitly passing the
goal into every trauma update call, substantially reduces the failure
mode where irrelevant named entities crowd out critical facts before
the token cap enforces consolidation.

\paragraph{Echidna trigger analysis.}
% TODO: fill in checkpoint frequency numbers from verbose logs
Echidna checkpoints more aggressively than necessary in document
ingestion settings, treating every topic shift between documents as a
compression trigger. In multi-hop retrieval, topic shifts between
documents are expected and not signals of reasoning state change. Tuning
Echidna's system prompt to distinguish document ingestion from free-form
conversation is a direct avenue for improvement.


% =========================================================
\section{Conclusion}
% =========================================================

We presented ReCompress, a two-component system addressing single-turn
and multi-turn LLM compression. The single-shot component demonstrates
that query-aware rewriting provides substantial gains on multi-hop QA
while emitting far fewer tokens than deletion, and that this capability
distills into a 1.5B model that transfers to a near-in-distribution benchmark
(directionally, but not yet significantly, to more dissimilar ones). The
multi-turn component demonstrates that explicit protection of critical facts
through tiered memory keeps solver context flat at no measurable accuracy cost.
We report the cost honestly: counting per-turn compression overhead, the naive
LLM-trigger version is more expensive in total tokens at short horizons. But we
also show the dominant overhead --- a per-turn Echidna LLM call that was
$\sim$98\% ``checkpoint'' anyway --- is removable for free, after which
RbD-Compress beats an \emph{uncached} growing-history agent on total tokens from
$\sim$6 turns onward (Section~\ref{sec:echidna-ablation}); against a
\emph{cached} deployment, flat context wins only when the context window is the
binding constraint or the solver dominates compression cost.

Query-aware rewriting and deletion-based compression serve complementary
operating regimes. Rewriting is best suited to tasks where query
specificity matters and distractor passages can be identified; deletion
is best suited to high-throughput, query-independent settings where
speed and cacheability are primary. Future work includes replacing
Echidna's language model call with a fine-tuned classifier, evaluating
on LongBench V2 with context lengths up to two million tokens, extending
RbD-Compress to agentic settings, and scaling the distillation dataset
to push directional wins on MuSiQue and SQuAD into statistical
significance.

\section*{Acknowledgments}
Built at the UC Berkeley AI Hackathon 2026 for The Token Company
Compression Challenge. We thank The Token Company for providing
bear-1.1 API access and for organizing the challenge.


% =========================================================
\clearpage
\appendix
% =========================================================

\section{Prompt Designs}

\paragraph{Trauma Extractor.}
\textit{You extract and maintain critical facts that must never be lost.
Given a new message and the current trauma memory, identify new critical
facts relevant to the goal: named entities, bridge facts, numbers, dates,
the user's core goal. Return JSON only with keys \texttt{update} and
\texttt{trauma\_memory}. Keep trauma memory under 50 tokens. Never
duplicate existing facts. For multi-hop QA, always pin entity names and
answers to sub-questions.}

\paragraph{Echidna.}
\textit{You are Echidna, the Witch of Greed. You observe conversations
with perfect clarity and decide when knowledge must be crystallized. You
receive the current turn, trauma memory, and checkpoint summary. Return
JSON only with keys \texttt{action}, \texttt{revert\_to},
\texttt{reason}, and \texttt{urgency}. Trigger a checkpoint when a topic
shift occurs, a reasoning hop has been resolved, or the token budget is
near its limit. Trigger a revert when a contradiction is detected or the
reasoning chain has collapsed. Otherwise pass. Always read trauma memory
before deciding.}

\paragraph{Compressor.}
\textit{Compress the following conversation history into a structured
checkpoint. Never compress, paraphrase, or omit any fact in trauma
memory. Preserve all named entities, bridging facts, and reasoning chain
steps. Target under 150 tokens total.}


\section{Distillation Training Trajectory}
\label{app:distillation}

\begin{table}[t]
\caption{Distillation training configurations and results on HotpotQA
vs.\ bear-1.1. n.s. = not significant (CI includes zero).}
\centering
\begin{tabular}{lcccc}
\toprule
Version & Examples & LoRA rank & Epochs & $\Delta$ vs bear (F1) \\
\midrule
v1 & 261   & 16 & 3 & +0.063 (n.s.) \\
v2 & 2,500 & 64 & 6 & +0.071 (n.s., overfit) \\
v3 & 5,000 & 32 & 3 & \textbf{+0.252} ($\star$) \\
\bottomrule
\end{tabular}
\end{table}


\section{Re:Zero Analogy}

\begin{table}[t]
\caption{Mapping between Re:Zero narrative elements and RbD-Compress
components.}
\centering
\begin{tabular}{ll}
\toprule
Re:Zero element & RbD-Compress component \\
\midrule
Return by Death checkpoint      & Checkpoint stack entry \\
Subaru's knowledge across loops & Trauma memory -- survives all operations \\
Loop delta (new info this loop) & Delta -- current turn only \\
Echidna, Witch of Greed         & Echidna trigger model \\
Subaru choosing his checkpoint  & User or Echidna-initiated revert \\
Satella's constraint            & Token budget hard limit \\
\bottomrule
\end{tabular}
\end{table}


% =========================================================
\bibliographystyle{unsrt}
\bibliography{references}
% =========================================================

\end{document}